\chapter{Evaluation and Related Work}

The evaluation asks whether the language boundary is executable and whether it
preserves useful compiler facts. It does not claim peak kernel performance
against mature tensor compilers. The relevant question is source structure:
whether axes, recurrence dependencies, and derivative targets remain available
without wrapper-style refactoring.

\section{Executable Credibility}

The implementation runs small kernels, autodiff examples, recurrence suites,
stdlib tests, simulation examples, optimization examples, and model-shaped
indexed programs. Unit tests inspect compiler metadata for recurrence storage
choices and lazy derivative values. Example tests compare selected programs
against NumPy references. Larger demonstrations check that the same syntax and
backend machinery scales beyond toy expressions.

\section{Source Boundary Comparisons}

For matrix multiplication, \einlang{} replaces a contraction string with source
indices:

\begin{lstlisting}[style=einlang]
let C[i, j] = sum[k](A[i, k] * B[k, j]);
\end{lstlisting}

The source-indexed form gives the compiler three facts before lowering:
\texttt{i} and \texttt{j} survive into the output, while \texttt{k} is consumed
by the reduction. That is more than a spelling difference from an
\texttt{einsum} string. It lets the resolver assign DefIds to the axes, lets
range analysis infer domains from shaped reads, and lets diagnostics point to
ordinary source syntax when a contraction dimension is inconsistent.

For convolution, it avoids a split between patch extraction and contraction:

\begin{lstlisting}[style=pythoncode]
patches = sliding_window_view(x, (KH, KW), axis=(2, 3))
y = np.einsum("ncijrs,ocrs->noij", patches, w)
\end{lstlisting}

The NumPy version is a useful baseline because it is explicit about the two
operations most libraries require: first build a windowed view, then contract
that view. The implementation claim for \einlang{} is that the sliding offsets
and contracted axes can live in one indexed source binding. That gives the
compiler a chance to check offset bounds, infer output shape, and choose a
backend path without asking the programmer to expose a patch tensor as a
separate source artifact.

For recurrence, it avoids requiring the programmer to encode the source
relation as a carry contract before the compiler has inspected later reads. The
advantage in each case is not shorter syntax by itself; it is that the compiler
receives one source object containing the facts that later passes need.

\section{What the Current Artifact Shows}

The current artifact supports four claims:

\begin{itemize}
\item Source coherence: indexed definitions, recurrent clauses, and derivative
requests can coexist in one binding environment.
\item Compiler leverage: range, shape, type, recurrence, and autodiff passes
consume facts that are still visible in the IR.
\item Runtime execution: the NumPy backend executes lowered tensor programs,
recurrences, stdlib calls, and local derivative requests.
\item Conservative optimization: finite recurrence storage is used only when
the compiler can prove bounded history and bounded observation.
\end{itemize}

\section{Evaluation Obligations}

The next evaluation step is to turn several current qualitative checks into
larger generated suites. For recurrence storage, the generator should vary
lookback distance, observed tail length, whole-tensor escape, symbolic extent,
dynamic tail indices, and multidimensional recurrence geometry. For autodiff,
the generator should compare scalar gradients, projected Jacobian entries,
rows, columns, and full materializations against numerical or framework
references where appropriate. For tensor lowering, the generator should compare
vectorized, scalar, hybrid, and call-scalar paths on the same mathematical
program to ensure that backend strategy changes do not change results.

\begin{lstlisting}[style=pythoncode,caption={Generated evaluation-suite outline.}]
generate_recurrence_suite():
    for lookback in [1, 2, 4]:
        for tail_observation in ["final", "last_k", "whole", "dynamic"]:
            program = synthesize_recurrence(lookback, tail_observation)
            result = compile_and_run(program)
            reference = run_numpy_reference(program.parameters)
            assert_close(result.value, reference)
            if tail_observation is bounded:
                assert compiler_storage(program).preserve_steps is finite
            else:
                assert compiler_storage(program).requires_full_output

generate_autodiff_suite():
    for shape_y, shape_x in scalar_and_tensor_shapes:
        program = synthesize_function_and_query(shape_y, shape_x)
        for projection in ["entry", "row", "column", "full"]:
            got = compile_and_run(project(program, projection))
            expected = numerical_or_framework_reference(program, projection)
            assert_close(got, expected)

generate_tensor_lowering_suite():
    for strategy in ["vectorized", "scalar", "hybrid", "call-scalar"]:
        program = synthesize_equivalent_tensor_program(strategy)
        got = compile_and_run(program)
        expected = direct_numpy_reference(program)
        assert_close(got, expected)
\end{lstlisting}

This generator would evaluate the compiler's smart paths directly: not only
whether a program returns the right value, but whether the compiler chose and
recorded the expected range, shape, recurrence-storage, autodiff, and
vectorization facts.

The generated-suite sketch is intentionally larger than a value-only test. A
recurrence case is not finished when the numerical output matches; the test
also has to assert whether the compiler proved a finite window or fell back to
full materialization for the right reason. An autodiff case is not finished
when a scalar gradient matches; it should also cover entry, row, column, and
full Jacobian projections so lazy lowering decisions are exercised. A tensor
lowering case is not finished when one strategy works; equivalent programs
should force vectorized, scalar, hybrid, and call-scalar paths.

This is the evaluation counterpart of the thesis implementation story. The
"smart" parts of the compiler are metadata-producing analyses, not only value
transformations. A strong generated suite therefore checks both outputs and the
facts attached to the IR. If a future optimization changes a backend strategy,
the suite should detect semantic regressions and also detect cases where the
compiler silently stopped proving the intended property.

The recurrence portion should also assert negative cases. A whole-tensor
observation, a dynamic tail index, or an unknown offset should force full
materialization even when the numerical result would match a circular-buffer
run for a particular input size. These tests protect the proof boundary: the
compiler is allowed to be conservative, but it is not allowed to claim a finite
window without a static reason. The generated suite should therefore check both
\texttt{preserve\_steps} values and fallback markers in the compiler metadata.

The autodiff portion needs the same two-layer structure. Numerical agreement
shows that the derivative value is correct for a sampled program, while
projection checks show that lazy lowering selected the intended row, column, or
entry path. A full materialization test is still useful, but it should not be
the only case. The thesis claim is specifically about local derivative requests
remaining shaped and demand-driven, so the evaluation should exercise uses that
ask for less than the dense Jacobian.

For tensor lowering, equivalent mathematical programs should be shaped so that
the compiler is forced onto different execution strategies. One version may be
fully vectorizable, another may include a loop-dependent function call, another
may include guards, and another may include nested reductions. Comparing all of
them against direct NumPy references checks semantic preservation, while
inspecting execution facts checks that the strategy classifier is still doing
the smart work the thesis describes.

\section{Limitations}

The strongest current result is language-design evidence backed by a concrete
compiler and runtime. The executable recurrence storage optimization targets
static offsets and one inferred finite recurrence axis. General
multidimensional mixed finite/full storage, dynamic offset reasoning,
automatic line-buffer inference, KV-cache inference, GPU code generation, full
whole-language type soundness, and broader IREE lowering remain future work.

\section{Related Work}

The closest neighboring ecosystems for everyday numerical programming are
NumPy, TensorFlow, PyTorch, and JAX. These systems provide the baseline for the
comparison here \cite{numpy2020,tensorflow2016,pytorch2019,jax2018}. They are
more mature and more general ecosystems than \einlang{}. The distinction is the
programming boundary: \einlang{} makes indexed definitions, recurrences, and
derivative requests native source constructs of one checked language rather
than separate host, library, and transformation interfaces.

TVM and Halide preserve compiler-facing tensor structure and expose powerful
scheduling or lowering interfaces \cite{tvm2018,halide2013}. \einlang{} sits
earlier in the source pipeline: it is concerned with how the programmer writes
the tensor relation, recurrence relation, and local derivative query before a
schedule is selected.

Automatic differentiation systems such as Zygote and DiffTaichi show that
differentiation can be integrated deeply into language or compiler structure
\cite{zygote2019,difftaichi2020}. \einlang{} shares the goal of treating
differentiation as a semantic concern, but its claim is narrower: derivative
requests are local expressions over named indexed and recurrent bindings.

The storage-window analysis connects to classical data-flow analysis, loop
transformations, locality optimization, software pipelining, and polyhedral
dependence analysis \cite{allen2001optimizing,muchnick1997advanced,lam1988software,rau1994iterative,feautrier1991data,bondhugula2008pluto,bastoul2004code}. The
specific contribution here is source placement: recurrent dependencies and
later observations are represented directly in the indexed program, allowing
storage planning before the computation has been lowered into backend loops.
